\documentclass[10pt,a4paper]{amsart}
\usepackage[margin=19mm]{geometry}
\usepackage{amsmath,amssymb,mathtools}
\usepackage[scr=rsfs,cal=euler]{mathalfa}
\usepackage{xfrac}
\usepackage[unicode,hidelinks,bookmarks=false]{hyperref}
\newcommand{\ie}{i.e.\ }
\newcommand{\cf}{cf.\ }
\newcommand{\resp}{respectively\ }
\newcommand{\Subsection}{\subsection}
\providecommand{\nobreakdash}{\nobreak\hbox{-}}
\setlength{\emergencystretch}{2em}
\multlinegap0pt
\usepackage[shortlabels]{enumitem}
\usepackage{mathtools}
\usepackage{amssymb}
\usepackage{xcolor}
\usepackage{tikz}
\newtheorem{lemma}{Lemma}
\newtheorem{theorem}{Theorem}
\title{PL Recognition After Two $S^{2}\times S^{2}$ Stabilisations is PSPACE-Hard}
\author{Rhuaidi Antonio Burke}
\date{Source: arXiv:2609.15890v1 (CC BY 4.0)}
\begin{document}
\maketitle

\noindent\emph{Source and licence.} PL Recognition After Two $S^{2}\times S^{2}$ Stabilisations is PSPACE-Hard. Version arXiv:2609.15890v1 (CC BY 4.0), distributed by arXiv under the Creative Commons Attribution 4.0 International licence, http://creativecommons.org/licenses/by/4.0/, as recorded in the arXiv licence metadata for this exact version and shown on the abstract page https://arxiv.org/abs/2609.15890v1. Full licence notice: \texttt{licenses/arxiv-2609.15890-CC-BY-4.0.txt}.\par\smallskip

\noindent\textit{Editorial scope.} Counted unit: the article's theorem thm:miller-tancer (source0.tex lines 501--522). The article's complete original proof of it is delivered with it (source0.tex lines 524--615, one \texttt{proof} environment of 92 lines). No mathematics is written, repaired or re-derived: the statement and the proof are the article's own lines, in the article's own notation, and no retained line is substituted or re-flowed.  Retained uncounted as the source-local context the proof needs: lines 452--467.  Nothing was omitted from any retained span, so the package carries no omission record.  No retained line contains a citation, so the wrapper carries no bibliography and no bibliography entry.  No retained line contains a figure environment, an image-inclusion command or drawing code, so no image is delivered and the package declares no asset dependency.  Editorial formatting only, in addition: an AMS \texttt{amsart} preamble that restates the article's own declarations and macros used by the retained lines, namely the environments \texttt{lemma}, \texttt{theorem} on separate counters, together with the standard packages those lines need.  The retained environments print in this document as Lemma 1, Theorem 1 in order of appearance, whereas the article prints the same items with the numbering of its own full document.  The complete native source of the article as distributed in the arXiv e-print for this exact version is bundled unchanged as \texttt{sources/210462/source0.tex}.
	\begin{lemma}
		\label{lem:branch-normal-form}
		Let $A$ be a group and $H=A*\langle t\rangle$. Let
		$D\subset\mathbb Z$ be finite. For each $d\in D$, choose exactly one
		element $z_d$ of one of the two forms
		\[
		z_d=t^{-d}a_dt^d,
		\qquad a_d\in A\text{ of infinite order}, \tag{P$_d$}
		\]
		or
		\[
		z_d=t^{-d}b_d^{-1}tb_dt^d,
		\qquad b_d\in A\setminus\{1\}. \tag{Q$_d$}
		\]
		Then $\{z_d:d\in D\}$ is a free basis for the subgroup it generates in $H$.
	\end{lemma}

	\begin{theorem}
		\label{thm:miller-tancer}
		For arbitrary $P,w,S$, the following hold.
		\begin{enumerate}[label=\textup{(\roman*)}]
			\item If $w=1$ in $G(P)$, then the map induced by the old generators is
			an isomorphism
			\[
			G(P^+)\xrightarrow{\ \cong\ }G(M(P,w,S)).
			\]
			\item If $w\ne1$ in $G(P)$, then the natural homomorphism
			\[
			G(P)\longrightarrow G(M(P,w,S))
			\]
			is injective.
		\end{enumerate}
		Consequently, if $P^+$ presents the trivial group, then
		\[
		M(P,w,S)\text{ presents the trivial group}
		\quad\Longleftrightarrow\quad
		w=1\text{ in }G(P).
		\]
	\end{theorem}

	\begin{proof}
		Suppose first that $w=1$ in $G(P)$. Then
		$[w,\beta]=1$ after imposing the relators of $P$, so (MT3) gives
		$\gamma^{-3}\beta\gamma^3=1$ and hence $\beta=1$. Relation (MT1)
		then gives $\gamma=1$, relation (MT2) gives $\alpha=1$, and
		$(\mathrm{MT4})_i$ gives $s_i=1$ for every $i$. Hence the old
		generators induce a homomorphism $G(P^+)\to G(M(P,w,S))$. Conversely,
		setting $\alpha=\beta=\gamma=1$ defines a homomorphism
		$G(M(P,w,S))\to G(P^+)$. The two maps are mutually inverse, proving (i).
		
		Now suppose that $w\ne1$ in $G(P)$. Put
		\[
		G_L=\bigl(G(P)*\langle\beta_L\rangle\bigr)*\langle\alpha\rangle,
		\qquad
		G_R=\langle\beta_R\rangle*\langle\gamma\rangle .
		\]
		In $G_L$, define
		\[
		\begin{aligned}
			u_0&=\beta_L,\\
			u_1&=\alpha^{-1}\beta_L\alpha,\\
			u_2&=\alpha^{-2}\beta_L^{-1}\alpha\beta_L\alpha^2,\\
			u_3&=\alpha^{-3}[w,\beta_L]\alpha^3,\\
			u_{3+i}&=\alpha^{-(i+3)}s_i\beta_L\alpha^{i+3},
			\qquad i=1,\dots,m'.
		\end{aligned}
		\]
		Apply Lemma~\ref{lem:branch-normal-form} with
		\[
		A=G(P)*\langle\beta_L\rangle,\qquad t=\alpha,
		\]
		and with the indices
		\[
		D=\{0,1,\ldots,m'+3\}.
		\]
		The elements $u_0,u_1,u_2,u_3$, and $u_{3+i}$ correspond respectively to the distinct indices $0,1,2,3$, and $i+3$. The element $u_2$ has form (Q$_2$), with $b_2=\beta_L$, while all the other displayed elements have form~(P$_d$) for their corresponding indices $d$.
		
		The element $\beta_L$ has infinite order. Since $w\ne1$, $[w,\beta_L]=w^{-1}\beta_L^{-1}w\beta_L$
		is cyclically reduced of syllable length four in
		$G(P)*\langle\beta_L\rangle$,
		and therefore has infinite order. 
		
		For each $i$, if $s_i=1$ in $G(P)$, then $s_i\beta_L=\beta_L$; otherwise, $s_i\beta_L$ is cyclically reduced of syllable length two. Thus $s_i\beta_L$ has infinite order in either case. All the hypotheses of Lemma~\ref{lem:branch-normal-form} are therefore satisfied, and it follows that
		\[
		u_0,u_1,\ldots,u_{m'+3}
		\]
		freely generate a free subgroup $A_L\le G_L$.
		
		In $G_R$, define
		\[
		\begin{aligned}
			v_0&=\beta_R,\\
			v_1&=\gamma^{-1}\beta_R^{-1}\gamma\beta_R\gamma,\\
			v_2&=\gamma^{-2}\beta_R^{-1}\gamma\beta_R\gamma^2,\\
			v_3&=\gamma^{-3}\beta_R\gamma^3,\\
			v_{3+i}&=\gamma^{-(i+3)}\beta_R\gamma^{i+3},
			\qquad i=1,\ldots,m'.
		\end{aligned}
		\]
		Apply Lemma~\ref{lem:branch-normal-form} with
		\[
		A=\langle\beta_R\rangle,\qquad t=\gamma,
		\]
		and the same index set $D=\{0,1,\ldots,m'+3\}$. The displayed elements have forms
		\[
		(\mathrm{P}_0),\ (\mathrm{Q}_1),\ (\mathrm{Q}_2),\ (\mathrm{P}_3),\ (\mathrm{P}_{i+3}),
		\]
		respectively, and their associated indices are pairwise distinct. Every coefficient required in Lemma~\ref{lem:branch-normal-form} is nontrivial and, where required, has infinite order. Hence
		$
		v_0,v_1,\ldots,v_{m'+3}
		$
		freely generate a free subgroup $A_R\le G_R$.
		
		The basis correspondence $u_j\mapsto v_j$ defines an isomorphism
		\(
		\phi\colon A_L\longrightarrow A_R.
		\)
		After the relation $u_0=v_0$ identifies $\beta_L$ with $\beta_R$, the
		presentation of the amalgamated free product
		$
		G_L*_{\phi}G_R
		$
		is precisely $M(P,w,S)$. The normal-form theorem for amalgamated free
		products gives injections of $G_L$ and $G_R$ into this group. Since
		$G(P)$ is a free factor of $G_L$, the natural map
		$
		G(P)\hookrightarrow G(M(P,w,S))
		$
		is injective, proving (ii). If $P^+$ presents the trivial group, (i) gives the yes-case,
		whereas (ii) gives a nontrivial subgroup in every no-case. This proves the
		final equivalence.
	\end{proof}

\end{document}
